\documentclass[11pt]{article}

% ---------- Page layout ----------
\usepackage[
    letterpaper,
    left=0.9in,
    right=0.9in,
    bottom=0.9in,
    top=0.5in
]{geometry}

% ---------- Fonts ----------
\usepackage[T1]{fontenc}
\usepackage[utf8]{inputenc}
\usepackage{lmodern}
\usepackage{microtype}
\usepackage{hyperref}
\usepackage{pgfplots}
\pgfplotsset{compat=1.18}
% ---------- Mathematics ----------
\usepackage{amsmath}
\usepackage{amssymb}
\usepackage{amsfonts}
\usepackage{mathtools}

% ---------- Figures ----------
\usepackage{graphicx}
\usepackage{float}
\usepackage[
    font=small,
    labelfont=bf
]{caption}
\usepackage{subcaption}

% ---------- Columbia logo ----------
\usepackage{tikz}
\usepackage{eso-pic}

% Put Columbia logo in upper-left corner of first page
\AddToShipoutPictureBG*{%
    \begin{tikzpicture}[remember picture,overlay]
        \node[
            anchor=north west,
            inner sep=0pt,
            opacity=0.75
        ]
        at ([xshift=0.45cm,yshift=-0.45cm]current page.north west)
        {\includegraphics[width=6cm]{mafn.png}};
    \end{tikzpicture}
}

% ---------- Slightly tighter figure spacing ----------
\setlength{\textfloatsep}{12pt}
\setlength{\floatsep}{10pt}
\setlength{\intextsep}{10pt}


% Python code formatting
\usepackage{listings}
\usepackage{xcolor}

% Dark code theme
\definecolor{codebg}{HTML}{15191F}
\definecolor{codefg}{HTML}{E6EDF3}
\definecolor{codekeyword}{HTML}{C792EA}
\definecolor{codestring}{HTML}{A5D6A7}
\definecolor{codecomment}{HTML}{7F8C98}
\definecolor{codenumber}{HTML}{8B949E}
\definecolor{codeborder}{HTML}{30363D}

\lstdefinestyle{darkpython}{
    language=Python,
    backgroundcolor=\color{codebg},
    basicstyle=\ttfamily\scriptsize\color{codefg},
    keywordstyle=\color{codekeyword}\bfseries,
    stringstyle=\color{codestring},
    commentstyle=\color{codecomment}\itshape,
    numberstyle=\tiny\color{codenumber},
    numbers=left,
    numbersep=12pt,
    stepnumber=1,
    frame=single,
    rulecolor=\color{codeborder},
    framesep=7pt,
    xleftmargin=15pt,
    xrightmargin=5pt,
    framexleftmargin=10pt,
    breaklines=true,
    breakatwhitespace=false,
    columns=fullflexible,
    keepspaces=true,
    showstringspaces=false,
    tabsize=4,
    captionpos=b,
    aboveskip=12pt,
    belowskip=12pt
}

% ---------- Document ----------
\begin{document}

\vspace*{0.25in}

\begin{center}
    {\Large\textbf{Nathaniel Coulter}\footnote{\href{https://github.com/Nathaniel-Coulter/Timeseries-StatInf-columbia}{github.com/Nathaniel-Coulter/Timeseries-StatInf-columbia}}}\\[3pt]
    {\large STAT 4263/5263 — Statistical Inference and Time Series Modeling}\\
    {\normalsize Homework 4 - 13 November 2026}
\end{center}

\vspace{0.5em}


\subsection*{Problem 1(a)}

For the ARMA(1,1) model
\[
X_t=0.7X_{t-1}+W_t+0.4W_{t-1},
\]
we have \(\phi=0.7\), \(\theta=0.4\), and
\(\sigma_W^2=1\).

Using the given autocovariance formulas,
\[
\begin{aligned}
\gamma(0)
&=1+\frac{(0.4+0.7)^2}{1-0.7^2}\\
&=3.372549,
\\[4pt]
\gamma(1)
&=0.4+0.7+
\frac{(0.4+0.7)^2(0.7)}{1-0.7^2}\\
&=2.760784.
\end{aligned}
\]

Therefore,
\[
\begin{aligned}
\rho(1)&=\frac{\gamma(1)}{\gamma(0)}
=0.818605,\\
\rho(2)&=0.7\rho(1)=0.573023,\\
\rho(3)&=0.7\rho(2)=0.401116.
\end{aligned}
\]

\textbf{Step 1:} For \(n=1\),
\[
\boxed{\phi_{11}=\rho(1)=0.818605}.
\]

The normalized prediction-error variance is
\[
v_1=1-\phi_{11}^2=0.329886.
\]

\textbf{Step 2:} For \(n=2\),
\[
\begin{aligned}
\phi_{22}
&=\frac{\rho(2)-\phi_{11}\rho(1)}
{1-\phi_{11}\rho(1)}\\
&=\frac{0.573023-(0.818605)^2}
{1-(0.818605)^2}\\
&=\boxed{-0.294314}.
\end{aligned}
\]

Updating the previous coefficient,
\[
\begin{aligned}
\phi_{21}
&=\phi_{11}-\phi_{22}\phi_{11}\\
&=0.818605-(-0.294314)(0.818605)\\
&=\boxed{1.059532}.
\end{aligned}
\]

The normalized prediction-error variance becomes
\[
\begin{aligned}
v_2
&=v_1(1-\phi_{22}^2)\\
&=(0.329886)(1-(-0.294314)^2)\\
&=0.301311.
\end{aligned}
\]

\textbf{Step 3:} For \(n=3\),
\[
\begin{aligned}
\phi_{33}
&=\frac{\rho(3)-\phi_{21}\rho(2)
-\phi_{22}\rho(1)}{v_2}\\
&=\frac{0.034907}{0.301311}\\
&=\boxed{0.115850}.
\end{aligned}
\]

Updating the remaining coefficients,
\[
\begin{aligned}
\phi_{31}
&=\phi_{21}-\phi_{33}\phi_{22}\\
&=1.059532-(0.115850)(-0.294314)\\
&=\boxed{1.093628},
\\[6pt]
\phi_{32}
&=\phi_{22}-\phi_{33}\phi_{21}\\
&=-0.294314-(0.115850)(1.059532)\\
&=\boxed{-0.417062}.
\end{aligned}
\]

Thus, the best linear predictor of \(X_4\)
using \(X_1,X_2,X_3\) is
\[
\boxed{
X_4^{(3)}
=1.093628X_3-0.417062X_2+0.115850X_1.
}
\]

\subsection*{Problem 1(b)}

The AI's mistake occurs in its calculation of
\(\phi_{33}\). It uses autocorrelations in the
numerator but the unnormalized prediction-error
variance in the denominator.

Specifically, the AI uses
\[
\gamma(0)(1-\phi_{11}^2)(1-\phi_{22}^2)
=1.0162.
\]

However, when the numerator is expressed
using autocorrelations, the denominator must
be the \textbf{normalized} prediction-error
variance:
\[
\begin{aligned}
v_2
&=(1-\phi_{11}^2)(1-\phi_{22}^2)\\
&=\boxed{0.301311}.
\end{aligned}
\]

Consequently,
\[
\boxed{
\phi_{33}
=\frac{0.034907}{0.301311}
=0.115850
}
\]
rather than \(0.0344\).

This error changes \(\phi_{33}\) and the
updated coefficients \(\phi_{31}\) and
\(\phi_{32}\), while leaving \(\phi_{11}\),
\(\phi_{21}\), and \(\phi_{22}\) unchanged.

It also invalidates the AI's conclusion
that an AR(2) model is sufficient.

\subsection*{Problem 1(c)}

The AI's conclusion that an AR(2) model
is sufficient is \textbf{incorrect}.

For an AR(2) process, the theoretical PACF
cuts off after lag 2, meaning
\[
\phi_{hh}=0,\qquad h>2.
\]

Using the approximate 95\% confidence bounds
with \(n=1440\),
\[
\pm\frac{1.96}{\sqrt{1440}}
=\boxed{\pm0.05165}.
\]

However, the correctly calculated
third partial autocorrelation is
\[
\boxed{
|\phi_{33}|=0.11585>0.05165.
}
\]

Thus, the lag-three PACF is outside these
bounds, contradicting the AI's argument
that the PACF cuts off after lag 2.

Furthermore, the underlying process is
ARMA(1,1), whose theoretical PACF generally
decays rather than cutting off at a finite lag.

Therefore, an AR(2) model is not an adequate
exact representation of the specified process.
If an AR approximation is required, higher-order
models should be considered and evaluated using
residual diagnostics.

The confidence bounds are an approximate
sampling guideline; since these PACF values
were calculated from the fitted theoretical
ACF rather than a directly estimated sample
PACF, they should not be interpreted as a
formal significance test.


\subsection*{Problem 2}

Let \(X_t\) be a stationary AR(1) process:
\[
X_t-\mu=\phi(X_{t-1}-\mu)+W_t,
\qquad |\phi|<1.
\]

For an AR(1), the autocovariance function is
\[
\gamma(h)=\gamma(0)\phi^{|h|}.
\]

We seek the best linear predictor of \(X_2\)
given \(X_1\) and \(X_3\):
\[
\hat X_2=a+b_1X_1+b_3X_3.
\]

The covariance matrix of \((X_1,X_3)\) is
\[
\Gamma=\gamma(0)
\begin{pmatrix}
1 & \phi^2\\
\phi^2 & 1
\end{pmatrix},
\]
and the covariance vector with \(X_2\) is
\[
\boldsymbol{\gamma}
=\gamma(0)
\begin{pmatrix}
\phi\\
\phi
\end{pmatrix}.
\]

The optimal coefficients satisfy
\[
\begin{pmatrix}
b_1\\b_3
\end{pmatrix}
=\Gamma^{-1}\boldsymbol{\gamma}.
\]

By symmetry, \(b_1=b_3=b\). The normal
equations give
\[
b(1+\phi^2)=\phi,
\]
so
\[
\boxed{b_1=b_3=\frac{\phi}{1+\phi^2}}.
\]

Therefore, the best linear predictor is
\[
\boxed{
\hat X_2
=\mu+\frac{\phi}{1+\phi^2}
\left[(X_1-\mu)+(X_3-\mu)\right].
}
\]

For a zero-mean AR(1), this simplifies to
\[
\boxed{
\hat X_2=\frac{\phi}{1+\phi^2}(X_1+X_3).
}
\]


\subsection*{Problem 4(a)}

The file \texttt{ibmclose.csv} contains 369 daily
closing prices. The original series, sample ACF,
and sample PACF are displayed below.

\begin{center}
\includegraphics[width=0.95\linewidth]{01_original_acf_pacf.png}
\end{center}

The series exhibits substantial changes in its
level, including a pronounced decline in the
later observations. The sample ACF decreases
very slowly and remains strongly positive over
many lags, while the PACF has a dominant spike
at lag one.

These features suggest that the original price
series is \textbf{nonstationary}. The
augmented Dickey--Fuller test gives
\(p\approx0.777\), so the unit-root null
hypothesis is not rejected.

\subsection*{Problem 4(b)}

We consider the three transformations
\[
g_1(X_t)=\log X_t,\qquad
g_2(X_t)=\sqrt{X_t},\qquad
g_3(X_t)=X_t^2,
\]
and calculate their first differences:
\[
Y_t^{(j)}=g_j(X_t)-g_j(X_{t-1}).
\]

\begin{center}
\includegraphics[width=0.95\linewidth]{02_transformed_differences.png}
\end{center}

Variance stability was compared using the
coefficient of variation of the standard
deviations across six consecutive time blocks.

\begin{center}
\begin{tabular}{lc}
\hline
Transformation & CV of block SDs\\
\hline
\(\Delta\log X_t\) & 0.5545\\
\(\Delta\sqrt{X_t}\) & 0.4570\\
\(\Delta X_t^2\) & 0.2291\\
\hline
\end{tabular}
\end{center}

Among these transformations, the squared-price
difference has the most stable variation by
this measure. We therefore select
\[
\boxed{Y_t=X_t^2-X_{t-1}^2}.
\]

For numerical stability during estimation,
the differences were rescaled by \(10^{-4}\).
The corresponding fitted series is
\[
Z_t=10^{-4}Y_t.
\]

For each transformation, ARMA(\(p,q\))
models were fitted with a mean, considering
\(p,q\in\{0,1,2,3\}\). The selected orders are

\begin{center}
\begin{tabular}{lccc}
\hline
Transformation & AIC & BIC & AICc\\
\hline
Logarithm & (3,2) & (0,0) & (3,2)\\
Square root & (3,2) & (0,0) & (3,2)\\
Square & (0,1) & (0,1) & (0,1)\\
\hline
\end{tabular}
\end{center}

For the selected squared-price difference,
the ARMA(0,1) model with a mean has
\[
\boxed{
\begin{aligned}
\mathrm{AIC}&=709.3898,\\
\mathrm{BIC}&=721.1140,\\
\mathrm{AICc}&=709.4557.
\end{aligned}
}
\]

Its maximum-likelihood parameter estimates
on the scaled series are
\[
\boxed{
\hat\mu=-0.022815,\qquad
\hat\theta_1=0.148306,\qquad
\hat\sigma_W^2=0.395910.
}
\]

Thus, the fitted model is
\[
Z_t=-0.022815+W_t+0.148306W_{t-1},
\]
where \(W_t\) has estimated variance
\(0.395910\).

\subsection*{Problem 4(c)}

For the AIC-selected ARMA(0,1) model, we
compare estimation with and without a mean.

The estimated mean with an intercept is
\[
\hat\mu=-0.022815.
\]

Its approximate 95\% confidence interval is
\[
\boxed{(-0.100216,\ 0.054587)}.
\]

Since this interval contains zero, there
is insufficient evidence that the mean
of the differenced series is nonzero.

The information criteria are

\begin{center}
\begin{tabular}{lccc}
\hline
Specification & AIC & BIC & AICc\\
\hline
With mean & 709.3898 & 721.1140 & 709.4557\\
Without mean & 707.7565 & 715.5727 & 707.7894\\
\hline
\end{tabular}
\end{center}

All three criteria favor excluding the mean.
Therefore, the preferred specification is
\[
\boxed{Z_t=W_t+0.149133W_{t-1}},
\]
with
\[
\boxed{\hat\sigma_W^2=0.396303}.
\]

\subsection*{Problem 4(d)}

We examine the standardized residuals, their
sample ACF, and the normal Q--Q plot for the
selected ARMA(0,1) model without a mean.

\begin{center}
\includegraphics[width=0.95\linewidth]{03_square_diagnostics.png}
\end{center}

The residual ACF is relatively small at
most lags, with no obvious persistent
autocorrelation pattern.

The Ljung--Box test results are
\[
\begin{array}{lc}
\hline
\text{Lag} & p\text{-value}\\
\hline
10 & 0.5157\\
20 & 0.1190\\
\hline
\end{array}
\]

Both p-values exceed \(0.05\), so the
null hypothesis of no residual
autocorrelation through the corresponding
lags is not rejected.

The Q--Q plot is approximately linear
in its central portion, but deviations
appear in the tails. Four standardized
residuals exceed three in absolute value.

The Jarque--Bera test also rejects normality,
with \(p<0.001\), indicating that the
Gaussian error assumption is imperfect.

Overall, the no-mean ARMA(0,1) model
provides an adequate parsimonious description
of the serial dependence in the transformed
differences, although the residual tails
suggest possible outliers or departures
from normality.


\subsection*{Problem 3(a)}

We generate 1,000 independent realizations
of the MA(2) process
\[
X_t=W_t+0.8W_{t-1}-0.1W_{t-2},
\qquad W_t\sim WN(0,1).
\]

Each simulated series contains \(n=200\)
observations, following a burn-in of
300 observations. A fixed random seed
is used for reproducibility.

For each series, MA(\(q\)) models are
estimated by Gaussian maximum likelihood
for \(q=0,1,2,3,4\), including a mean.

Model selection is based on
\[
\begin{aligned}
\mathrm{AIC}&=-2\ell+2k,\\
\mathrm{BIC}&=-2\ell+k\log n,\\
\mathrm{AICc}&=\mathrm{AIC}
+\frac{2k(k+1)}{n-k-1},
\end{aligned}
\]
where \(\ell\) is the maximized
log-likelihood and \(k=q+2\) counts
the mean, MA coefficients, and
innovation variance.

The estimated selection probabilities are

\begin{center}
\begin{tabular}{cccc}
\hline
\(q\) & AIC & BIC & AICc\\
\hline
0 & 0.000 & 0.000 & 0.000\\
1 & 0.429 & 0.778 & 0.445\\
2 & 0.401 & 0.212 & 0.404\\
3 & 0.099 & 0.008 & 0.094\\
4 & 0.071 & 0.002 & 0.057\\
\hline
\end{tabular}
\end{center}

The true order is \(q=2\). AIC and
AICc select the correct order with
probabilities 0.401 and 0.404,
respectively, while BIC selects
the correct order with probability 0.212.

BIC favors the simpler MA(1) model
in 77.8\% of simulations. This likely
reflects the relatively small magnitude
of the second MA coefficient,
\(\theta_2=-0.1\), combined with
BIC's stronger complexity penalty.

AIC selects orders above the true
order in 17.0\% of simulations,
compared with 15.1\% for AICc.
Thus, AICc slightly reduces
overfitting relative to AIC.

\subsection*{Problem 3(b)}

We generate 1,000 independent realizations
of the AR(3) process
\[
X_t=0.8X_{t-1}-0.2X_{t-2}
-0.1X_{t-3}+W_t,
\]
where \(W_t\sim WN(0,1)\).

Each series contains \(n=200\) observations,
following a burn-in of 300 observations.

For each simulated series, AR(\(p\))
models with a mean are fitted by Gaussian
maximum likelihood, considering
\(p=0,1,2,3,4,5\).

The same information criteria are used,
with \(k=p+2\).

The estimated selection probabilities are

\begin{center}
\begin{tabular}{cccc}
\hline
\(p\) & AIC & BIC & AICc\\
\hline
0 & 0.000 & 0.000 & 0.000\\
1 & 0.005 & 0.047 & 0.006\\
2 & 0.421 & 0.753 & 0.435\\
3 & 0.405 & 0.184 & 0.405\\
4 & 0.107 & 0.016 & 0.100\\
5 & 0.062 & 0.000 & 0.054\\
\hline
\end{tabular}
\end{center}

The true order is \(p=3\). Both AIC
and AICc select the correct order
with probability 0.405, while BIC
selects it with probability 0.184.

BIC selects AR(2) in 75.3\% of
simulations, again tending to
underestimate the true order.
The relatively small third AR
coefficient, \(\phi_3=-0.1\),
makes distinguishing AR(3) from
AR(2) difficult at this sample size.

AIC selects orders above the true
order in 16.9\% of simulations,
compared with 15.4\% for AICc.
BIC rarely overfits but frequently
selects an underspecified model.

\textbf{Conclusion.}
Across both experiments, AIC and AICc
recover the true order more frequently
than BIC at \(n=200\), although neither
consistently identifies the correct model.

BIC favors simpler models because of
its larger penalty for additional
parameters. AICc moderates the
overfitting tendency of AIC through
its finite-sample correction.

These findings illustrate the trade-off
between underfitting and overfitting
in finite samples. BIC's asymptotic
consistency does not guarantee the
highest probability of selecting the
true model at a particular finite
sample size.




\subsection*{Problem 5(a)}

The given AR(2) model is
\[
X_t-1.2X_{t-1}+0.5X_{t-2}=W_t,
\qquad W_t\sim WN(0,1).
\]

Its characteristic polynomial is
\[
\Phi(z)=1-1.2z+0.5z^2.
\]

Solving \(\Phi(z)=0\) gives
\[
z_{1,2}=1.2\pm i\sqrt{0.56}.
\]

Both roots satisfy
\[
|z_{1,2}|=\sqrt{2}>1.
\]

Therefore, the process is \textbf{causal}
and admits a convergent MA(\(\infty\))
representation
\[
X_t=\sum_{j=0}^{\infty}\psi_jW_{t-j}.
\]

Matching coefficients in
\[
(1-1.2B+0.5B^2)
\sum_{j=0}^{\infty}\psi_jB^j=1,
\]
we obtain
\[
\psi_0=1,\qquad
\psi_1=1.2,
\]
and the recursion
\[
\psi_j=1.2\psi_{j-1}-0.5\psi_{j-2}.
\]

Thus,
\[
\begin{aligned}
\psi_0&=\boxed{1},\\
\psi_1&=\boxed{1.2},\\
\psi_2&=1.2(1.2)-0.5(1)
=\boxed{0.94},\\
\psi_3&=1.2(0.94)-0.5(1.2)
=\boxed{0.528}.
\end{aligned}
\]

An unexpected one-megawatt shock corresponds
to an innovation of magnitude \(W_t=1\).
Since observations are 15 minutes apart,
the predicted additional power use is
\[
\boxed{
\begin{aligned}
15\text{ minutes: }&1.2\text{ MW},\\
30\text{ minutes: }&0.94\text{ MW},\\
45\text{ minutes: }&0.528\text{ MW}.
\end{aligned}
}
\]

These values represent changes relative
to the forecast without the shock.

\subsection*{Problem 5(b)}

For a stationary AR(2) process, the
Yule--Walker equations give
\[
\rho(1)=1.2-0.5\rho(1),
\]
and hence
\[
\rho(1)=0.8.
\]

For higher lags,
\[
\rho(h)=1.2\rho(h-1)-0.5\rho(h-2).
\]

Therefore,
\[
\begin{aligned}
\rho(2)&=1.2(0.8)-0.5=0.46,\\
\rho(3)&=1.2(0.46)-0.5(0.8)=0.152.
\end{aligned}
\]

Using the Durbin--Levinson recursion,
\[
\phi_{11}=\rho(1)=0.8,
\]
\[
\phi_{22}
=\frac{\rho(2)-\rho(1)^2}
{1-\rho(1)^2}
=\frac{0.46-0.64}{1-0.64}
=-0.5,
\]
and
\[
\phi_{21}
=\phi_{11}-\phi_{22}\phi_{11}
=1.2.
\]

At lag three,
\[
\begin{aligned}
\phi_{33}
&=\frac{\rho(3)-\phi_{21}\rho(2)
-\phi_{22}\rho(1)}
{1-\phi_{21}\rho(1)-\phi_{22}\rho(2)}\\
&=\frac{0.152-1.2(0.46)+0.5(0.8)}
{1-1.2(0.8)+0.5(0.46)}\\
&=\frac{0}{0.27}.
\end{aligned}
\]

Thus,
\[
\boxed{\phi_{33}=0}.
\]

\subsection*{Problem 5(c)}

The AI's three answers contain the
following errors.

\begin{enumerate}
\item \textbf{Incorrect causality condition.}
The AI applies the AR(1) condition
\(|\phi_1|<1\) to an AR(2) model.
For an AR(2), causality requires both
roots of the characteristic polynomial
to lie outside the unit circle.
Here their moduli are \(\sqrt{2}>1\),
so the model is causal.

\item \textbf{Incorrect impulse-response
recursion.}
The AI uses a positive coefficient on
\(\psi_{j-2}\), but the correct recursion is
\[
\psi_j=1.2\psi_{j-1}-0.5\psi_{j-2}.
\]
Consequently, the correct coefficients are
\[
(\psi_0,\psi_1,\psi_2,\psi_3)
=(1,1.2,0.94,0.528).
\]
The AI's growing coefficients arise from
the incorrect recursion and do not
demonstrate noncausality.

\item \textbf{Confusing ACF with PACF.}
The value \(0.152\) is the lag-three
autocorrelation \(\rho(3)\), not the
partial autocorrelation \(\phi_{33}\).
The correct value is
\[
\boxed{\phi_{33}=0}.
\]
The theoretical PACF of an AR(2) process
cuts off after lag two, so an AR(3)
model is not required by this evidence.
\end{enumerate}

The AI's impulse-response calculation
does not independently confirm its
causality claim because it uses the same
incorrectly specified recursion.
A computation based on an incorrect
recurrence cannot serve as independent
evidence for the original model.

The correct characteristic roots and
convergent impulse-response coefficients
support the causal AR(2) specification.


\section*{AI-use disclosure}
OpenAI ChatGPT was used to draft/check the mathematical derivations, evaluate the provided AI response in Problem 2, and generate parts of the code that made the plots in Problem 3, it also suggested a few tweaks to my original code for cleanliness that I obliged. 

\appendix



\end{document}